\BourbakiExercisesSection

\begin{enumerate}
\def\labelenumi{\arabic{enumi}.}
\tightlist
\item
  Let E be a set, \(A \subset E \times E\) and \((x, y) \mapsto x \top y\) , \((x, y) \in A\) , an internal law of composition not everywhere defined on E. Given two subsets X and Y of E, let \(X \top Y\) denote the set of elements of E of the form \(x \top y\) with \(x \in X\) , \(y \in Y\) and \((x, y) \in A\) . A law of composition is thus defined which is everywhere defined on \(\mathfrak{P}(E)\) .
\end{enumerate}

If \((\mathbf{X}_{\alpha})_{\alpha \in \mathbf{A}}\) and \((\mathbf{Y}_{\beta})_{\beta \in B}\) are any two families of subsets of E, then

\[
\left(\bigcup_ {\alpha \in A} X _ {\alpha}\right) \top \left(\bigcup_ {\beta \in B} Y _ {\beta}\right) = \bigcup_ {(\alpha , \beta) \in A \times B} (X _ {\alpha} \top Y _ {\beta}).
\]

\begin{enumerate}
\def\labelenumi{\arabic{enumi}.}
\setcounter{enumi}{1}
\item
  Let \(\top\) be a law of composition not everywhere defined on a set E. Let E' be the subset of \(\mathfrak{P}(E)\) consisting of the sets \(\{x\}\) where \(x\) runs through E and the empty subset \(\varnothing\) of E; show that E' is a stable subset of \(\mathfrak{P}(E)\) under the law \((X,Y)\mapsto X\top Y\) (Exercise 1); deduce that, if \(\overline{E}\) denotes the set obtained by adjoining (Set Theory, II, § 4, no. 8) to E an element \(\omega\) , the law \(\top\) can be extended to \(\overline{E} \times \overline{E}\) , so that the law \(\top\) is identical with the law induced on E by this extended law.
\item
  Let \(\top\) be a law of composition not everywhere defined on E.
\end{enumerate}

\begin{enumerate}
\def\labelenumi{(\alph{enumi})}
\item
  For the law \((\mathbf{X}, \mathbf{Y}) \mapsto \mathbf{X} \top \mathbf{Y}\) (Exercise 1) between subsets of \(\mathbf{E}\) to be associative, it is necessary and sufficient that, for all \(x \in \mathbb{E}\) , \(y \in \mathbb{E}\) , \(z \in \mathbb{E}\) , if one of the two sides of the formula \((x \top y) \top z = x \top (y \top z)\) is defined, so is the other and is equal to it (use Exercise 1 to show that the condition is sufficient).
\item
  If this condition holds, show that Theorem 1 of no. 3 can be generalized as follows: if one of the two sides of formula (3) is defined, the other is defined and is equal to it.
\end{enumerate}

\begin{enumerate}
\def\labelenumi{\arabic{enumi}.}
\setcounter{enumi}{3}
\tightlist
\item
  \begin{enumerate}
  \def\labelenumii{(\alph{enumii})}
  \tightlist
  \item
    Given a set E, let \(\Phi\) be the set of mappings into E of any subset of E; if f, g, h are three elements of \(\Phi\) , show that if the composition \((f \circ g) \circ h\) is defined so is \(f \circ (g \circ h)\) , but that the converse is not true; if these two compositions are defined they are equal.
  \end{enumerate}
\end{enumerate}

\begin{enumerate}
\def\labelenumi{(\alph{enumi})}
\setcounter{enumi}{1}
\tightlist
\item
  Let \(\mathfrak{F}\) be a family of non-empty subsets of \(E\) such that no two have an element in common and let \(\Psi\) be the subset of \(\Phi\) consisting of the bijective mappings of a set of \(\mathfrak{F}\) onto a set of \(\mathfrak{F}\) . Show that, under the law induced by the law \(f \circ g\) on \(\Psi\) , the condition of Exercise 3(a) holds.
\end{enumerate}

\begin{enumerate}
\def\labelenumi{\arabic{enumi}.}
\setcounter{enumi}{4}
\item
  Show that the only triples \((m, n, p)\) of natural numbers \(\neq 0\) such that \((m^n)^p = m^{n^p}\) are: \((1, n, p)\) , \(n\) and \(p\) being arbitrary; \((m, n, 1)\) , \(m\) and \(n\) arbitrary and \((m, 2, 2)\) where \(m\) is arbitrary.
\item
  Let \(\top\) be a law of composition on a set E; let A be the subset of E consisting of the elements \(x\) such that \(x \top (y \top z) = (x \top y) \top z\) for all \(y \in E\) , \(z \in E\) ; show that A is a stable subset of E and that the law induced on A by \(\top\) is associative.
\item
  If \(\top\) is an associative law on \(\mathbf{E}\) and \(a\) and \(b\) two elements of \(\mathbf{E}\) , the sets \(\{a\} \top \mathbf{E}, \mathbf{E} \top \{b\}, \{a\} \top \mathbf{E} \top \{b\}, \mathbf{E} \top \{a\} \top \mathbf{E}\) are stable subsets of \(\mathbf{E}\) under the law \(\top\) .
\item
  Let \(\top\) be an associative law on E and \(a\) an element of E; for all \(x \in \mathrm{E}\) , \(y \in \mathrm{E}\) , we write \(x \perp y = x \top a \top y\) : show that the law \(\bot\) is associative.
\item
  On a set E the mappings \((x,y)\mapsto x\) and \((x,y)\mapsto y\) are opposite associative laws of composition.
\item
  Let X and Y be any two subsets of a set E and let \(X \top Y = X \cup Y\) if \(X \cap Y = \varnothing\) , \(X \top Y = E\) if \(X \cap Y \neq \varnothing\) ; show that the law of composition thus defined on \(\mathfrak{P}(E)\) is associative and commutative.
\item
  Let \(\top\) be an associative law on \(\mathbf{E}\) and \(\mathbf{A}\) and \(\mathbf{B}\) two subsets of \(\mathbf{E}\) which are stable under this law; show that, if \(\mathbf{B} \top \mathbf{A} \subset \mathbf{A} \top \mathbf{B}\) , \(\mathbf{A} \top \mathbf{B}\) is a stable subset of \(\mathbf{E}\) .
\item
  The only distinct natural numbers \(\neq 0\) which are permutable under the law \((x,y)\mapsto x^{y}\) are 2 and 4.
\item
  Show that under the law \((\mathbf{X}, \mathbf{Y}) \mapsto \mathbf{X} \circ \mathbf{Y}\) between subsets of \(E \times E\) the centre is the set consisting of \(\varnothing\) and the diagonal.
\item
  Show that under the law of composition \(f \circ g\) between mappings of E into E the centre consists of the identity mapping.
\item
  A law \(\top\) on a set \(E\) is called idempotent if all the elements of \(E\) are idempotent (no. 4) under this law, that is if \(x \top x = x\) for all \(x \in E\) . Show that, if a law \(\top\) on \(E\) is associative, commutative and idempotent, the relation \(x \top y = y\) is an order relation on \(E\) ; if we write \(x \leqslant y\) , any two elements \(x, y\) of \(E\) admit a least upper bound (under this order relation) equal to \(x \top y\) . Obtain the converse.
\item
  Let \(\top\) be an associative and commutative law of composition on a set E. Let \(n\) be an integer \(>0\) . The mapping \(x \mapsto \frac{n}{\top} x\) of E into E is a morphism.
\end{enumerate}

